\chapter{The tools}
\label{ch:tools}

\section{The goal}

The data was fixed; the instrument was not. Before any experiment could run, one
practical question had to be settled: \emph{given one free graphics card and a
few weeks, which model, which training method and which software are actually
feasible?}

This toolchain was decided in the first week and never changed. What changed over
the internship was the data and the training signal, never the instrument, and
that is what makes the later comparisons between runs meaningful.

\section{The tools, and what each one does}

% F5a -- The tools, as a stack: what sits on top of what, and who does what.
\begin{figure}[!htbp]
\centering
\fitwidth{%
\begin{tikzpicture}[font=\small,
  lay/.style={draw=clrule,rounded corners=2pt,align=center,inner sep=6pt,
              minimum height=1.35cm,text width=11.4cm,anchor=center},
  box/.style={draw=clrule,rounded corners=2pt,align=center,inner sep=6pt,
              minimum height=1.75cm,text width=3.35cm,font=\footnotesize,anchor=center},
  wide/.style={box,text width=5.35cm},
  side/.style={draw=clrule,fill=white,rounded corners=2pt,align=center,
               inner sep=6pt,text width=3.3cm,font=\footnotesize,anchor=center},
  tag/.style={font=\scriptsize,text=clrule,anchor=east},
]
% --- layer 4: my own code -------------------------------------------------
\node[lay,fill=clmodel!12,draw=clmodel] (mine) at (0,4.75)
  {\textbf{My notebook}\\[2pt]\footnotesize
   the custom training loop: which mistakes count more, which set of numbers
   answers for which kind of user, and the whole grading};
\node[tag] at (-6.1,4.75) {written by me};

% --- layer 3: the three libraries ----------------------------------------
\node[box,fill=clsoft] (tr)   at (-3.95,2.55)
  {\textbf{Transformers}\\[2pt] loads the model and runs a standard training
   loop};
\node[box,fill=clsoft] (peft) at (0,2.55)
  {\textbf{PEFT}\\[2pt] freezes the model and adds the small trainable part};
\node[box,fill=clsoft] (trl)  at (3.95,2.55)
  {\textbf{TRL}\\[2pt] the format the training file is written in};
\node[tag] at (-6.1,2.55) {the libraries};

% --- layer 2: model + 4-bit storage --------------------------------------
\node[wide,fill=clink!14] (qwen) at (-2.95,0.35)
  {\textbf{Qwen2.5-0.5B}\\[2pt] the model itself, 500 million numbers,
   downloaded already trained on general text};
\node[wide,fill=clink!14] (bnb) at (2.95,0.35)
  {\textbf{bitsandbytes}\\[2pt] stores those frozen numbers in 4 bits each
   instead of 16, so they fit};
\node[tag] at (-6.1,0.35) {what is trained};

% --- layer 1: the machine -------------------------------------------------
\node[lay,fill=clnonsig!35] (colab) at (0,-1.65)
  {\textbf{Google Colab, free tier: one NVIDIA T4 graphics card, 16\,GB}\\[2pt]
   \footnotesize rented in the browser, a few hours a day, and liable to
   disconnect};
\node[tag] at (-6.1,-1.65) {where it runs};

% --- the side column: analysis -------------------------------------------
\node[side,anchor=west] (an) at (6.55,1.45)
  {\textbf{pandas}\\ \textbf{scikit-learn}\\ \textbf{scipy}\\[4pt]
   preparing the files, splitting people into train and test, and checking the
   split is fair.\\[3pt] No modelling here.};

\draw[clrule,dashed] (6.15,-2.45) -- (6.15,5.55);
\end{tikzpicture}}
\caption{The tools, and the job of each one}
\label{fig:toolchain}

\end{figure}


\section{How much of the model to change}
\label{sec:tools-peft}

A language model is a stack of very large tables of
numbers.\footnote{\textbf{Weights}, or \textbf{parameters}: the numbers inside
the model that turn an input into an output. ``Frozen'' means left exactly as its
original authors trained them.} Adapting one somebody else has already
trained\footnote{\textbf{Pre-trained}: already trained, by someone else, on a very
large body of general text. Fine-tuning starts from those numbers rather than
from nothing, which is what makes it affordable.} can be done at three very
different costs, and what separates them is not accuracy. It is memory.

% F5b -- Full fine-tuning, LoRA and QLoRA, drawn side by side.
\begin{figure}[!htbp]
\centering
\fitwidth{%
\begin{tikzpicture}[font=\small,
  big/.style={draw=clrule,fill=clnonsig!30,minimum width=2.5cm,minimum height=2.5cm},
  small/.style={draw=clgain,fill=clgain!25,minimum width=2.5cm,minimum height=0.42cm},
  smallv/.style={draw=clgain,fill=clgain!25,minimum width=0.42cm,minimum height=2.5cm},
  q/.style={draw=clmodel,fill=clmodel!18,minimum width=2.5cm,minimum height=2.5cm},
  lbl/.style={font=\scriptsize,text=clrule,align=center},
]
% ---------------- (a) full fine-tuning ----------------
\node[big,fill=clloss!25,draw=clloss] (fa) at (0,0) {};
\node[lbl,above=3pt of fa] {\textbf{(a) Full fine-tuning}};
\node[lbl,below=3pt of fa,align=center] {every weight is updated\\ \textcolor{clloss}{100\,\% trainable}\\ does not fit on a free GPU};
\node[font=\scriptsize,text=clloss] at (fa) {$W$};

% ---------------- (b) LoRA ----------------
\node[big] (fb) at (4.6,0) {};
\node[font=\scriptsize,text=clrule] at (fb) {$W$ \emph{frozen}};
\node[smallv] (bA) at (6.35,0.55) {};
\node[small]  (bB) at (6.35,-1.0) {};
\node[font=\scriptsize] at (bA) {$A$};
\node[font=\scriptsize] at (bB) {$B$};
\node[lbl,above=3pt of fb] {\textbf{(b) LoRA}};
\node[lbl,below=17pt of fb,align=center,xshift=0.9cm]
  {the frozen matrix is left alone; only the two\\ thin \textcolor{clgain}{$A$} and
   \textcolor{clgain}{$B$} matrices are trained\\ \textcolor{clgain}{$\approx$2\,\% trainable}};

% ---------------- (c) QLoRA ----------------
\node[q] (fc) at (10.4,0) {};
\node[font=\scriptsize,text=clmodel,align=center] at (fc) {$W$ \emph{frozen}\\ \emph{and 4-bit}};
\node[smallv] (cA) at (12.15,0.55) {};
\node[small]  (cB) at (12.15,-1.0) {};
\node[font=\scriptsize] at (cA) {$A$};
\node[font=\scriptsize] at (cB) {$B$};
\node[lbl,above=3pt of fc] {\textbf{(c) QLoRA}, used here};
\node[lbl,below=17pt of fc,align=center,xshift=0.9cm]
  {same, plus the frozen matrix is stored\\ at 4 bits per number instead of 16\\
   \textcolor{clmodel}{$\approx$1/4 of the frozen memory}};
\end{tikzpicture}}

\caption{What LoRA and QLoRA do}
\label{fig:lora}
\figtext{%
\figpara{The grey square is the model.}{A language model is essentially a stack
of very large tables of numbers, and adapting one to a new task means changing
numbers in those tables. The three panels are the three ways of doing it, drawn
at the same scale.}

\figpara{(a) Change everything.}{Every number in the table is modified. This has
the highest ceiling in principle, but while the run lasts each number has to be
held in the graphics card's memory together with about three more used to keep
track of how to adjust it. The whole thing is several times larger than the
model, and it does not fit on a free card.}

\figpara{(b) LoRA: leave the table alone and add two thin ones beside
it.}{\cite{hu2022lora} The big table is frozen and never touched. Two narrow
tables are added next to it, and only those are trained. Two narrow tables
standing in for a change to a wide one is the whole of the idea, and here it
means about 2 numbers in every 100 are ever allowed to move.}

\figpara{(c) QLoRA: the same, with the frozen table compressed.}{\cite{dettmers2023qlora}
This is the method used throughout the report. On top of LoRA, the frozen table
is stored with four bits per number instead of sixteen, which quarters the memory
it occupies. The compression is safe precisely because those numbers are frozen:
they are read, never updated. The added tables stay at full precision.}

\figpara{Why this decides the shape of the whole internship.}{Not accuracy, but
waiting time. With (c), one training run fits in half an hour on a free graphics
card, so nine weeks buy thirty experiments instead of two.}}

\end{figure}


I chose QLoRA, and I chose it on hardware grounds rather than on expected
quality: it was the only one of the three under which my first candidate model
could be loaded at all on a free graphics card. That choice held for the rest of
the internship, and it is what made the volume of experiments possible. It is
also what made the per-group routing of Chapter~\ref{ch:groups} cheap enough to
try at all.

\section{Which model, and why the small one won}
\label{sec:tools-model}

The first real run used \textbf{Mistral-7B} \cite{jiang2023mistral}, a
general-purpose model with seven billion numbers in it.\footnote{``7B'' means
seven billion numbers. Qwen2.5-0.5B has five hundred million, fourteen times
fewer. The count is the model's capacity, and also how much memory it occupies.}
A deliberately small first experiment, about thirty minutes over 1.2\,\% of the
training data, produced 16.87\,\% correct next-antenna answers, which was
encouraging as a signal that the pipeline worked end to end. A Llama model tried
at the same size did not even return well-formed antenna names.

Scaled up to a full run on 400 training people, however, Mistral-7B spent roughly
five and a half hours on the free card to reach a best score of 41.1\,\% and then
stopped improving. The same data given to \textbf{Qwen2.5-0.5B-Instruct}
\cite{qwen2025}, fourteen times smaller, ran forty-nine passes over the data in
about seventy-two minutes, and twenty-nine minutes once the fixes of
Chapter~\ref{ch:first} were in. Mistral-7B was abandoned.

The reasoning, stated at the time and still my view, is that the task is narrow.
Naming the next antenna out of 369 requires no knowledge of the world, no
reasoning in English and no ability to follow instructions. It requires learning
how one small geographical area is moved through. A seven-billion-number model
brings capacity the task cannot use and charges for it in waiting time, by a
factor of roughly twenty-three per run, which over nine weeks is the difference
between thirty experiments and two.

\section{How a day is written down for the model}

The model cannot read a day; it reads symbols.

% F3b -- One training example: the JSONL record, and the same line as tokens.
\begin{figure}[!htbp]
\centering

{\textbf{\panelnum{fig:training-example}{a} One person's day, as it is stored on disk.} One line of the
training file, spread over several lines here so that it can be read. Nothing is
a conversation in any real sense: the three labels are simply the slots the
training library expects to find.\par}
\vspace{3pt}

\begin{lstlisting}
{"conversations": [
  {"role": "system",    "content": "You are an AI that generates a user's
                                    cell-visit trajectory."},
  {"role": "user",      "content": "Here is the sequence for this user."},
  {"role": "assistant", "content": "User 2832452 | 200 events |
                                    BSOCHO1 BSOCHO1 UKVCHE102 BSOCHO1 ..."}
]}
\end{lstlisting}

\vspace{2pt}
{\footnotesize
\begin{tabular}{@{}p{2.2cm}p{10.75cm}@{}}
\texttt{system} & What the task is. The same words on every line of the file. \\
\texttt{user}   & The request. Also the same everywhere: a placeholder, not a
                  question. \\
\texttt{assistant} & The person's actual day. This is what the model has to
                  learn to produce, and the only part its mistakes are counted
                  on. \\
\end{tabular}\par}

\vspace{0.55cm}
\hrule height 0.4pt
\vspace{0.45cm}

{\textbf{\panelnum{fig:training-example}{b} The same line, once it has been cut into symbols.} This is what
the model actually reads. Each antenna name is \emph{one} symbol, not a spelling;
each event is preceded by its hour. Only the dark blue positions are graded.\par}
\vspace{5pt}

\fitwidth{%
\begin{tikzpicture}[font=\scriptsize,node distance=1.2pt,
  tok/.style={draw=clrule,minimum height=0.55cm,inner xsep=3.5pt,inner ysep=0pt},
  hr/.style={tok,draw=clmove,fill=clmove!20},
  cl/.style={tok,draw=clmodel,fill=clmodel!18},
  pfx/.style={tok,fill=clsoft,text=clrule},
]
\node[pfx] (t1) {\textless system\textgreater};
\node[pfx,right=of t1] (t2) {\dots};
\node[pfx,right=of t2] (t3) {User 2832452};
\node[pfx,right=of t3] (t4) {\textbar};
\node[hr,right=of t4]  (t5) {\textless H07\textgreater};
\node[cl,right=of t5]  (t6) {BSOCHO1};
\node[hr,right=of t6]  (t7) {\textless H09\textgreater};
\node[cl,right=of t7]  (t8) {BSOCHO1};
\node[hr,right=of t8]  (t9) {\textless H12\textgreater};
\node[cl,right=of t9]  (t10) {UKVCHE102};
\node[hr,right=of t10] (t11) {\textless H18\textgreater};
\node[cl,right=of t11] (t12) {BSOCHO1};
\node[pfx,right=of t12] (t13) {\dots};

\draw[decorate,decoration={brace,amplitude=4pt,mirror},clrule]
  ($(t1.south west)+(0,-2pt)$) -- ($(t4.south east)+(0,-2pt)$)
  node[midway,below=5pt,font=\scriptsize,text=clrule,align=center] {opening\\ \emph{not graded}};
\draw[decorate,decoration={brace,amplitude=4pt,mirror},clmodel]
  ($(t5.south west)+(0,-2pt)$) -- ($(t13.south east)+(0,-2pt)$)
  node[midway,below=5pt,font=\scriptsize,text=clmodel,align=center]
  {the day itself: only the 4 blue positions are graded};
\end{tikzpicture}}

\vspace{0.15cm}
\vspace{-0.15cm}
\replegend{\swatch{clsoft}~opening: excluded from the grade \qquad
           \swatch{clmove!20}~hour: \emph{given}, never guessed \qquad
           \swatch{clmodel!18}~antenna name: the only positions graded}

\caption{One person's day, as the model receives it}
\label{fig:training-example}

\end{figure}


Three choices are visible in that picture, and each is a single sentence with a
large consequence.

\paragraph{One antenna, one symbol.} All 369 antenna names were added to the
model's vocabulary\footnote{\textbf{Vocabulary}: the fixed set of symbols a
language model can read and produce. \textbf{Tokenizer}: the component that cuts
an input into those symbols.} as symbols of their own. Left alone, the tokenizer
would have split \cid{BSOCHO1} into four meaningless fragments, and the model
would have had to spell an antenna correctly to be right, which is precisely the
failure the Llama attempt showed.

\paragraph{One hour, one symbol.} From the sixth iteration each event is preceded
by a symbol for its hour of day, \cid{<H00>} to \cid{<H23>}. Hours rather than
minutes or seconds means twenty-four learnable symbols instead of 86{,}400
unlearnable ones, and one symbol rather than a written-out time preserves the
history the model can see, since it may only look at 512 symbols at a
time.\footnote{\textbf{Context window}: the maximum number of symbols the model
may look at in one go. Not the same as the \emph{context fraction} of
Chapter~\ref{ch:first}, which is a property of the measurement.}

\paragraph{Only the antennas are graded.} The opening boilerplate and the hour
symbols are excluded from the score.\footnote{\textbf{Masking} a position means
excluding it from the score: the model still \emph{reads} the symbol but is
neither rewarded nor penalised for producing it.} The model is trained, and
graded, only where an antenna name has to be produced.

\section{How one training run works}

% F5d -- How one training run actually works, step by step.
\begin{figure}[!htbp]
\centering
\fitwidth{%
\begin{tikzpicture}[font=\small,
  st/.style={draw=clrule,fill=clsoft,rounded corners=3pt,align=center,
             inner sep=6pt,text width=2.85cm,minimum height=2.75cm},
  ar/.style={-{Stealth[length=5pt]},clrule,very thick},
  num/.style={circle,draw=clmodel,fill=white,inner sep=1.2pt,
              font=\scriptsize\bfseries,text=clmodel},
]
\node[st] (a) at (0,0)
  {{\footnotesize\bfseries\color{clmodel} Take a slice}\\[2pt]\footnotesize
   a stretch of one training person's day: hour, antenna, hour, antenna};
\node[st] (b) at (3.55,0)
  {{\footnotesize\bfseries\color{clmodel} Hide the next antenna}\\[2pt]\footnotesize
   the model reads the slice and names what it thinks comes next};
\node[st] (c) at (7.10,0)
  {{\footnotesize\bfseries\color{clmodel} Measure the surprise}\\[2pt]\footnotesize
   compare with the antenna really visited; a confident wrong answer costs more
   than a hesitant one};
\node[st] (d) at (10.65,0)
  {{\footnotesize\bfseries\color{clmodel} Nudge the small part}\\[2pt]\footnotesize
   only the added numbers move; the 500 million original ones never do};

\draw[ar] (a) -- (b); \draw[ar] (b) -- (c); \draw[ar] (c) -- (d);
\node[num] at (a.north west) {1}; \node[num] at (b.north west) {2};
\node[num] at (c.north west) {3}; \node[num] at (d.north west) {4};

% --- the loop back --------------------------------------------------------
\draw[ar] (d.south) -- ++(0,-0.65) -- ++(-10.65,0) -- (a.south);
\node[font=\scriptsize,text=clrule] at (5.32,-2.42)
  {repeat, slice after slice, until every training person has been read once:
   that is \textbf{one pass}};

% --- the checkpoint rule --------------------------------------------------
\node[draw=clgain,fill=clgain!10,rounded corners=3pt,align=center,inner sep=6pt,
      text width=13.4cm,font=\small,anchor=north] (ck) at (5.32,-2.85)
  {{\footnotesize\bfseries\color{clgain} 5. At the end of every pass: grade the
    model on people it is not training on, save a copy, and keep only the best
    copy of the whole run}\\[2pt]\footnotesize
   On this task the best copy arrives between the fourth and the ninth pass, and
   everything after it is worse. Keeping the last copy instead of the best one
   costs 15.5 correct answers per hundred (Chapter~\ref{ch:first}).};
\end{tikzpicture}}
\caption{How one training run works}
\label{fig:training-loop}
\figtext{%
\figpara{Steps 1 to 4 are the whole of training.}{A slice of somebody's day is
shown to the model with the next antenna hidden. The model names an antenna. The
answer is compared with the truth and turned into a single number measuring how
surprised the model was, so that being confidently wrong costs more than
hesitating. That number is then used to nudge the model's adjustable numbers
very slightly in the direction that would have made it less wrong. Then the next
slice.}

\figpara{Only a small part of the model ever moves.}{Step 4 is where the method
of Figure~\ref{fig:lora} bites. The 500 million numbers the model was published
with are frozen and never touched; the nudging applies only to a small set of
numbers added beside them, about 2 out of every 100. That is what makes a
training run fit in half an hour on a free graphics card.}

\figpara{Step 5 is not an optimisation, it is a safety rule.}{Left to itself, a
model trained on a few hundred short days stops learning the habits of the
population and starts learning the days themselves by heart. When that happens
its score on people it has never seen turns around and falls, while its score on
its own training data keeps improving, so the training figures alone will happily
report success for another forty passes. Grading on separate people after each
pass, and keeping the best copy rather than the last, is what prevents that, and
it was the largest single improvement of the whole internship.}}
\end{figure}


Two further mechanisms were added to this loop from the sixth iteration onward,
and later chapters depend on both.

\paragraph{Some mistakes are made to count more.} At step 3 the surprise measured
at a given position can be multiplied according to the hour that position falls
in. A mistake at seven in the evening under a multiplier of four costs four times
the same mistake at two in the afternoon. No data is duplicated; the optimiser is
simply told which mistakes matter.

\paragraph{Some slices are shown more often.} At step 1, slices containing an
interesting moment can be drawn three times more often than the rest. This is a
different lever: it changes how often the question is asked, not how much a wrong
answer costs.

Both turned out to help only where the targeted positions were genuinely the hard
ones, which is the subject of Chapters~\ref{ch:time} and~\ref{ch:groups}.

\section{The machine: one free graphics card}

Every run reported in this document was trained on \textbf{Google
Colab}\footnote{\textbf{Google Colab}: a service that runs a Python notebook on a
graphics card rented in the browser, free for a few hours a day and liable to
disconnect. \textbf{T4} is the entry-level NVIDIA card it hands out.} on an
NVIDIA \textbf{T4}. That card carries 16\,GB of memory, and that single figure is
the constraint behind every choice in this chapter. The daily free quota is why
the experiment schedule is one training per iteration rather than broad sweeps of
settings.

Three standard libraries do work that never touches the
model.\footnote{\textbf{pandas}, \textbf{scikit-learn} and \textbf{scipy}: the
three standard Python libraries for, respectively, manipulating tables of data,
machine-learning utilities, and statistical tests.} They produce the balanced
split between training and test people, check that the two halves really do have
comparable histories, and run the resampling test
\cite{efron1993bootstrap}\footnote{\textbf{Bootstrap}: error bars obtained by
re-scoring the models on many random redraws of the test people, with no
assumption about the shape of the underlying distribution.} that decides the
central question of Chapter~\ref{ch:scale}.

\section{What I reused from my colleagues}

Three pieces came from the team. The daily extracts and the selection of people
who move and stay in the area came out of Nino's own analysis tooling; he ran the
extraction for me whenever I needed a new regime, and those files are the data on
which most of this report's modelling rests. The move-or-stay framing, separating
the moments where a person really goes somewhere from the moments where nothing
happens, is Tiphain and Benjamin's, and it is the idea behind the filtering of
Chapter~\ref{ch:data} and behind reporting in-window and out-of-window scores
separately. The technique of levelling every person's record count, tested in
Chapter~\ref{ch:scale}, is Tiphain's.

% F5c -- The other branch: one prompt sent to a frozen model.
\begin{figure}[t]
\centering

{\footnotesize\textbf{One of the prompts of the prompting branch,} as it is sent
to a frozen model. Quoted from Tiphain and Benjamin's work and shortened at the
two ellipses; nothing else is altered.\par}
\vspace{3pt}

\begin{lstlisting}
User's day started at 00:00 at station VSOYRO. [...] They stayed at BSOROT
until 05:34, then left for station VOYQZ, arriving at 06:15. [...] It is now
07:45 and the user is at VSOYRO.
Predict the station the user will visit next. Rely only on their movement
pattern and the stations around them. Possible stations include: BKVOR4,
BKVOF2, BKVS2D, BSOROT, BSVOLOK, DRVOAN, DSXVOP, LKVTSE. Which station are
they most likely to go to next?
Look at the user's daily flow (loops, pauses, commuting, travel habits).
Explain your choice in a few natural sentences, then end with the station in
this format:
FINAL ANSWER: [STATION]
\end{lstlisting}

\caption{What one prompt looks like}
\label{fig:prompting}
\figtext{%
The same problem, asked the other way. Three differences with the training
record of Figure~\ref{fig:training-example} are worth seeing, because they are
what makes the two branches answer different questions. The day is narrated in
English, stop by stop and with clock times, so one event costs a whole clause
where the fine-tuned model reads two symbols, and a long user stops fitting in
the model's window long before a short one does. The candidate cells are handed
over inside the prompt, eight of them here, and the answer comes back as free
text that has to be parsed, which is what the \texttt{FINAL ANSWER} convention
is for; the fine-tuned model instead chooses among 369 symbols it already
carries in its vocabulary. And since no weight ever moves, nothing learned
about one user is available for the next: everything the model is to know about
this population has to be restated in every prompt, whereas the whole point of
the branch this report takes is that the population is what the weights end up
holding. Neither is a better instrument. Asking a general model what it already
knows about a day described in words, and asking what can be learned from the
population itself, are two questions, and the second is this report's.}
\end{figure}


\section{Looking back}
\label{sec:retro-tools}

One of these choices was solved by another. I adopted QLoRA because storing the
frozen numbers in four bits was the only way to load a seven-billion-number model
on a free card, and a week later I abandoned that model for one fourteen times
smaller, which would have fitted without the compression. Keeping QLoRA was
harmless, but it illustrates engineering around a constraint instead of removing
it: the cheaper answer was to make the model smaller, not the numbers coarser.

The one setting I never revisited is how much of the model I allowed to move.
About 2 numbers in every 100 changed in every run of the internship, from four
hundred training people to twenty-one thousand, and that fraction was never
scaled with the data. Chapter~\ref{ch:scale} is where the omission bites.
